
In various application domains such as networking or cyber-security, large volumes of time series datasets are continuously generated and collected that are valuable for monitoring systems performance. Accurately detecting anomalies within these datasets are essential for identifying networks faults, device/system malfunctions, security breaches or service degradations that can significantly impact operations and user quality of experience (QoE). However, labeling real-time series data for anomaly detection is a time-consuming and challenging task due to the large amount of available time-series data and the scarcity of labeled anomalies in production systems. 

As a result, unsupervised time series anomaly detection (AD) has received significant attention in machine learning research community and several papers have proposed different techniques for this purpose. 
The proposed AD techniques can be categorized into five groups: reconstruction-based \cite{audibert2020usad, park2018multimodal, su2019robust, xu2018unsupervised}, distance-based \cite{breunig2000lof}, one-class classification-based \cite{pmlr-v80-ruff18a, scholkopf1999support}, isolation-based \cite{liu2008isolation}, and generative-based \cite{Zenati2018AdversariallyLA, zhou2019beatgan}. These techniques aim to learn the patterns of normality exclusively from normal data and detect anomalies based on deviations from the learned normality. The deviations are computed using metrics such as reconstruction error or distance to normal centroids. However, these techniques have some inherent limitations. Firstly, they struggle to discriminate anomalies that are close to normal samples because they do not leverage information about the anomalous class which are not used during the training phase. Secondly, the presence of unknown anomalies in the training sets introduces the challenge of \textit{data contamination}, which can negatively impact the performance of these models.

To enhance the performance of time series AD, contrastive learning (CL) \cite{chopra2005learning}
emerges as a promising approach, increasingly applied to various classification and forecasting tasks \cite{Eldele2021TimeSeriesRL, franceschi2019unsupervised, Yue2021TS2VecTU, Zheng2023SimTSRC}. CL has gained popularity across multiple disciplines including computer vision and natural language processing (NLP). Recently, its application in time series analysis has attracted attention. The fundamental concept of CL is to learn data representations by contrasting positive views and negative views. Specifically, through data augmentation, input data is transformed into multiple views, which are then contrasted in the latent space. This contrastive process aims to cluster similar views while repelling dissimilar ones. Following the training phase, the encoder is employed for downstream tasks. The effectiveness of contrastive learning stems from its ability to learn transformation-invariant properties through data augmentations. 

However, in time series analysis tasks, the exploration of data augmentation techniques has not yet been as extensive as in the field of computer vision. Additionally, some previous works in this field have used CL for anomaly detection. Nevertheless, they did not utilize a similarity function specifically designed for time series data, resulting in an inefficient exploitation of the temporal aspect inherent in multivariate time series data, which are nevertheless crucial for modeling. To address these limitations, we propose in this chapter an innovative end-to-end method called contrastive learning for anomaly detection in time series (CATS). We evaluate this method using large and several time series datasets including cloud gaming and system monitoring applications. 
%\abdelkader{REMPLACE LA DERNIERE PHRASE: We evaluate this method using several time series datasets including cloud gaming and system monitoring applications.}
\contributions{
Specifically, the main contributions of this chapter can be summarized as follows:
%which aims to enhance time series anomaly detection through a novel temporal representation learning. Specifically, our contributions can be summarised as follows:
\begin{itemize}
    \item Using Dynamic Time Warping (DTW) similarity, we propose a novel DTW-based temporal contrastive learning loss to efficiently model multivariate time series .
    \item We employ negative data augmentations to generate synthetic anomalies to establish a realistic out-of-order distribution that contrasts with normal instances in the training set.
    \item We conduct extensive empirical experiments on large real-world and popular benchmark time series datasets. Furthermore, we conduct experiments on time series datasets in the networking domain, collected from cloud gaming platforms. We demonstrate the effectiveness of our proposed framework and its generalization capabilities compared to its counterparts. 
    %which are also utilized by state-of-the-art techniques. Furthermore, to 
    %demonstrate the effectiveness of our proposed framework and its generalization capabilities, we undertake comparative evaluation with its competitors using time series datasets collected from cloud gaming platforms.
    \item We conduct ablation studies to evaluate the effectiveness of each component of the proposed method, assess the impact of data augmentation, data contamination and the influence of hyperparameters.
\end{itemize}
}

% The remainder of the paper is organized as follows. Section \ref{chap_5:related_work} presents the related work. Section \ref{chap_5:method} describes CATS method for anomaly detection in time series and Section \ref{chap_5:evaluation} the experimental setup.
